\section{Background and Related Work}
\label{sec:related}

\paragraph{DBpedia.} DBpedia~\cite{auer2007dbpedia, lehmann2015dbpedia}
extracts several datasets from every Wikipedia article: labels, short and long
abstracts, page identifiers, images, categories, redirects, external links,
geographic coordinates, interlanguage links and infobox data. Infobox data is
published twice. The \emph{raw} extraction turns every template parameter
into a property of the generic \term{dbp:} namespace. The \emph{mapping-based}
extraction uses community-maintained mappings from templates to the DBpedia
Ontology (\term{dbo:}), a hand-built schema of classes and typed properties,
which yields cleaner data at lower coverage. Each resource receives an IRI
derived from its article title, such as
\term{http://dbpedia.org/resource/Hanoi}. Non-English chapters publish data
from their own Wikipedia edition under their own namespace and link their
resources to the English edition. Kontokostas et al.~\cite{kontokostas2012greek}
describe the Greek chapter and the issues that internationalisation raises:
IRIs in a non-Latin script, language-specific templates, and interlinking
across editions. Our project follows the same model for Vietnamese. We reuse
the dataset types of DBpedia and its \term{dbo:} vocabulary, and we add a small
domain ontology where DBpedia lacks the terms we need.

\paragraph{Linked Data and the five-star scheme.} The Linked Data principles
ask publishers to name things with HTTP URIs, to return useful information in
RDF when a URI is looked up, and to link to other
URIs~\cite{bernerslee2006linked, heath2011linked}. The five-star scheme grades
open data: one star for data on the web under an open licence, two for
machine-readable structure, three for non-proprietary formats, four for using
URIs to identify things, and five for linking to other data. Tasks~3 and~4 of
the assignment correspond to the fourth and fifth stars. We follow the W3C
recommendations for serving such data: \emph{303 See Other} redirects with
content negotiation, so that one URI identifies an entity and leads either to
an HTML page or to an RDF document~\cite{sauermann2008cooluris,
hyland2014bestpractices}, and a VoID description of the dataset and its
linksets~\cite{alexander2011void}. Data are expressed in RDF~1.1~\cite{cyganiak2014rdf}
and queried with SPARQL~1.1~\cite{harris2013sparql}.

\paragraph{Wikidata.} Wikidata~\cite{vrandecic2014wikidata} is a collaborative
knowledge base whose items carry typed statements with qualifiers, and
\emph{sitelinks} to the article about the same item in every Wikipedia
edition. The Wikidata Query Service exposes this data as a SPARQL
endpoint~\cite{malyshev2018wdqs}. We use Wikidata in three ways. It selects the
entities through their class and country statements. It supplies precise
dates and quantities, which are easy to parse in Wikidata but often written as
free text in Vietnamese infoboxes. And its English sitelinks give the
English DBpedia resource that corresponds to each entity.

\paragraph{Ontology engineering and reasoning.} We designed the ontology from
\emph{competency questions}: questions the knowledge base must be able to
answer, written before the schema and later used as
tests~\cite{gruninger1995methodology}. To derive DBpedia types and shortcut
relations from our own vocabulary, we use the OWL~2~RL
profile~\cite{motik2012owl2profiles}. Its rules can be applied by forward
chaining in time polynomial in the size of the data, and the derived triples
can be materialised and published alongside the asserted ones.

\paragraph{Question answering over knowledge graphs.} Retrieval-augmented
generation grounds the answer of a language model in retrieved
evidence~\cite{lewis2020rag}. Over a knowledge graph the retrieval step can be
a structured query: the model translates the question into SPARQL, the query
runs on the graph, and the model answers from the result rows. Our
question-answering tab follows this pattern. The model receives a summary of
the schema and a few example queries, and invalid or empty queries are sent
back to it for repair.
